\documentclass[11pt,a4paper]{article}
\usepackage[utf8]{inputenc}
\usepackage[T1]{fontenc}
\usepackage{lmodern}
\usepackage[french]{babel}
\providecommand{\og}{\guillemotleft\nobreakspace}
\AtBeginDocument{\renewcommand{\contentsname}{Table des mati\`eres}\renewcommand{\abstractname}{R\'esum\'e}\renewcommand{\figurename}{Figure}}
\providecommand{\fg}{\nobreakspace\guillemotright}
\usepackage[margin=2.4cm]{geometry}
\usepackage{microtype}
\usepackage{amsmath,amssymb}
\usepackage{booktabs,longtable,array,tabularx,multirow}
\usepackage{xcolor}
\usepackage{graphicx}
\usepackage{tikz}
\usetikzlibrary{arrows.meta,positioning,shapes.geometric,calc,fit}
\usepackage{listings}
\usepackage[most]{tcolorbox}
\usepackage{enumitem}
\usepackage{fancyhdr}
\usepackage{hyperref}
\usepackage{float}

\definecolor{accent}{RGB}{20,80,140}
\definecolor{okc}{RGB}{30,120,60}
\definecolor{warnc}{RGB}{170,90,0}
\definecolor{badc}{RGB}{160,30,30}
\definecolor{grayc}{RGB}{90,90,90}
\hypersetup{colorlinks=true,linkcolor=accent,urlcolor=accent,citecolor=accent,
  pdftitle={Rapport technique -- LiDAR-camera fusion, AEB, ADAS Validation Studio et agent LLM},
  pdfauthor={Mohamed Eladeb}}

\setlist{itemsep=2pt,topsep=4pt}
\setlength{\parindent}{0pt}
\setlength{\emergencystretch}{3em}
\setlength{\parskip}{5pt}
\renewcommand{\arraystretch}{1.15}
\pagestyle{fancy}\fancyhf{}
\fancyhead[L]{\small Fusion LiDAR--cam\'era, AEB et agent LLM sous CARLA}
\fancyhead[R]{\small Rapport technique}
\fancyfoot[C]{\thepage}
\renewcommand{\headrulewidth}{0.3pt}

\lstset{basicstyle=\ttfamily\footnotesize,breaklines=true,columns=fullflexible,frame=single,
  rulecolor=\color{grayc},backgroundcolor=\color{black!3},keepspaces=true,upquote=true,
  literate={é}{{\'e}}1 {è}{{\`e}}1 {à}{{\`a}}1 {ê}{{\^e}}1 {ç}{{\c c}}1 {ô}{{\^o}}1 {û}{{\^u}}1 {ï}{{\"i}}1 {î}{{\^i}}1 {É}{{\'E}}1 {→}{{$\rightarrow$}}1}

% Statuts (distinction exigée dans tout le projet)
\newcommand{\impl}{\textcolor{warnc}{\textbf{impl\'ement\'e}}}
\newcommand{\teste}{\textcolor{accent}{\textbf{test\'e}}}
\newcommand{\valide}{\textcolor{okc}{\textbf{valid\'e en simulation}}}
\newcommand{\plan}{\textcolor{badc}{\textbf{planifi\'e}}}
\newcommand{\code}[1]{\texttt{#1}}

\newtcolorbox{honnete}[1][Lecture honn\^ete]{colback=accent!5,colframe=accent,fonttitle=\bfseries\small,title=#1,boxrule=0.5pt,left=4pt,right=4pt,top=2pt,bottom=2pt}
\newtcolorbox{limite}[1][Limite]{colback=badc!5,colframe=badc,fonttitle=\bfseries\small,title=#1,boxrule=0.5pt,left=4pt,right=4pt,top=2pt,bottom=2pt}
\newtcolorbox{choix}[1]{colback=okc!5,colframe=okc,fonttitle=\bfseries\small,title={Choix : #1},boxrule=0.5pt,left=4pt,right=4pt,top=2pt,bottom=2pt}

\title{\Large\bfseries Rapport technique d\'etaill\'e\\[4pt]
\large Fusion tardive LiDAR--cam\'era, freinage d'urgence automatique (AEB) valid\'e en boucle ferm\'ee,\\
\emph{ADAS Validation Studio} et agent LLM de g\'en\'eration de sc\'enarios d'essai\\[4pt]
\normalsize Projet personnel sous CARLA 0.9.16}
\author{Mohamed Eladeb (ENSTA Paris, 3A parcours robotique)}
\date{7 novembre 2026}

\begin{document}
\maketitle
\thispagestyle{fancy}

\begin{abstract}
Ce document d\'ecrit l'ensemble du projet \code{carla-lidar-camera-fusion} : la cha\^ine de perception
(collecte synchronis\'ee, \'etiquetage automatique, d\'etecteur YOLO fine-tun\'e, estimation de distance par LiDAR),
le suivi de cibles par filtre de Kalman, le freinage d'urgence automatique valid\'e en boucle ferm\'ee dans CARLA,
l'interface web \emph{ADAS Validation Studio} (API FastAPI, campagnes d'essais) et l'agent qui convertit une demande
en langage naturel en sc\'enario JSON ex\'ecutable. Pour chaque brique sont donn\'es : le but, la m\'ethode, les choix
de conception et leur justification, les r\'esultats mesur\'es, les erreurs rencontr\'ees et corrig\'ees, et surtout
\emph{ce qui n'est pas d\'emontr\'e}. Le rapport distingue partout ce qui est \impl, \teste, \valide{} ou \plan.
\end{abstract}

\tableofcontents
\newpage

%=====================================================================
\section{Contexte, objectifs et r\`egles de travail}
%=====================================================================
\subsection{Objectif du projet}
Le projet est un travail personnel men\'e pour se rapprocher du domaine automobile embarqu\'e (ADAS, perception,
validation). Il est parti d'une question de perception --- \emph{que gagne-t-on \`a fusionner un LiDAR et une cam\'era ?}
--- puis a \'et\'e \'etendu en trois \'etapes successives, chacune construite sur la pr\'ec\'edente :
\begin{enumerate}
\item une cha\^ine de perception caméra + LiDAR (d\'etection, distance, rejet de fausses alarmes) ;
\item un AEB fond\'e sur cette perception, \'evalu\'e en boucle ferm\'ee dans le simulateur CARLA ;
\item une application web (\emph{ADAS Validation Studio}) pour lancer et analyser des essais et des campagnes, et un agent LLM qui
      traduit une demande textuelle (\og voiture \`a l'arr\^et \`a 50~m, j'arrive \`a 70~km/h \fg) en sc\'enario valide.
\end{enumerate}

\subsection{R\`egles de m\'ethode (cahier des charges de conduite du projet)}
Ces r\`egles ont guid\'e chaque d\'ecision et expliquent le ton prudent du rapport :
\begin{itemize}
\item \textbf{V\'erifier avant de supposer} : lire le code et les sorties r\'eelles plut\^ot que d'affirmer de m\'emoire.
\item \textbf{Ne jamais inventer de r\'esultat} : tout chiffre cit\'e provient d'une sortie de commande ou d'un fichier de \code{reference/}.
\item \textbf{Distinguer} \impl{} (le code existe), \teste{} (des tests automatiques passent), \valide{} (mesur\'e dans CARLA dans des
      conditions pr\'ecis\'ees) et \plan{} (pas encore fait).
\item \textbf{Expliquer les cons\'equences} d'un choix (ce qu'il am\'eliore, ce qu'il d\'egrade).
\item \textbf{Dire les limites d'acc\`es} : l'assistant n'a pas de terminal sur la machine ; les commandes sont lanc\'ees par l'auteur,
      qui colle les sorties, et fait lui-m\^eme les \emph{commits} Git.
\end{itemize}

\subsection{Environnement}
\begin{tabularx}{\textwidth}{lX}\toprule
Syst\`eme & Windows ; GPU NVIDIA RTX 4050 (6~Go) \\
Simulateur & CARLA 0.9.16 (mode synchrone, pas de 0,05~s = 20~Hz), serveur local \code{localhost:2000} \\
Python & 3.12, environnement virtuel \code{venv} ; \code{torch 2.11.0+cu128}, \code{ultralytics 8.4.155}, NumPy, OpenCV \\
Application & FastAPI + Uvicorn ; page HTML/JavaScript sans framework (le choix React/TypeScript initialement envisag\'e est report\'e) \\
Tests & \code{pytest}, sans CARLA \\
D\'epôt & \code{C:\textbackslash Projets\textbackslash carla-lidar-camera-fusion} (Git ; donn\'ees \code{data/} et \code{runs/} non versionn\'ees) \\
\bottomrule\end{tabularx}

%=====================================================================
\section{Vue d'ensemble et structure du projet}
%=====================================================================
\subsection{Cha\^ine de traitement}
\begin{figure}[H]
\centering
\resizebox{\textwidth}{!}{\begin{tikzpicture}[node distance=7mm and 6mm,
  box/.style={draw=accent,rounded corners,fill=accent!6,minimum height=9mm,minimum width=22mm,align=center,font=\footnotesize},
  arr/.style={-{Stealth[length=2mm]},thick,accent}]
\node[box] (cam) {Cam\'era RGB\\800$\times$600, FOV 90$^\circ$};

\node[box,right=14mm of cam] (yolo) {YOLO26n\\fine-tun\'e\\(1 classe)};
\node[box,right=of yolo] (fus) {Fusion tardive\\distance LiDAR\\dans la bo\^ite};
\node[box,above=of fus] (lid) {LiDAR 64 voies\\1,2\,M pts/s, 70~m};
\node[box,right=of fus] (trk) {Tracker\\Kalman\\(x, y, v, a)};
\node[box,right=of trk] (aeb) {AEB\\machine d'\'etats\\$a_{req}$};
\node[box,right=of aeb] (veh) {Frein\\de l'ego};
\draw[arr] (cam)--(yolo); \draw[arr] (yolo)--(fus); \draw[arr] (fus)--(trk); \draw[arr] (trk)--(aeb); \draw[arr] (aeb)--(veh);
\draw[arr] (lid.south) -- (fus.north);
\draw[arr,dashed,grayc] (veh.south) -- ++(0,-.9) -| node[pos=.25,below,font=\scriptsize,text=grayc]{boucle ferm\'ee (CARLA, 20~Hz)} (cam.south);
\end{tikzpicture}}
\caption{Cha\^ine de perception et de d\'ecision. En boucle ferm\'ee, le frein appliqu\'e change ce que les capteurs voient \`a l'image suivante.}
\end{figure}

\subsection{Arborescence}
\begin{lstlisting}
carla-lidar-camera-fusion/
  config / requirements.txt, README.md
  scripts/        00_rig_test ... 13_agent_eval        (une etape = un script)
  src/fusion/
    config.py      parametres du banc (capteurs, episodes, split)
    geometry.py    reperes, projection camera, transformations
    rig.py         capteurs CARLA, meteo, synchronisation, etiquetage automatique
    perception/    fusion camera-LiDAR (sans import carla : testable hors simulation)
    adas/          tracker de Kalman, AEB, banc cinematique, injection de defauts
    agent/         schema JSON, validation, prompt, client LLM, regles, metriques
  backend/
    app/main.py    API FastAPI (runs, campagnes, agent)
    worker.py      worker d'observation (YOLO + LiDAR) en sous-processus
    static/index.html   interface web
  reference/       resultats figes (CSV, jsonl, hash du modele)
  tests/           test_aeb, test_faults, test_perception, test_fusion_lateral, test_api, test_agent
  data/ runs/ models/  (non versionnes)
\end{lstlisting}

\subsection{Scripts, du plus ancien au plus r\'ecent}
\begin{longtable}{@{}>{\ttfamily}p{3.6cm}p{11.4cm}@{}}\toprule
\textrm{Script} & R\^ole \\\midrule\endhead
00\_rig\_test & test du banc de capteurs \\
01\_collect & collecte (\'episodes $\times$ m\'et\'eos clair/pluie/brouillard), reprise possible apr\`es arr\^et \\
02\_view, 03\_audit & visualisation et audit du jeu de donn\'ees \\
04\_make\_yolo & conversion au format YOLO (une classe \emph{vehicle}) \\
05\_eval\_yolo & \'evaluation du d\'etecteur par distance et par m\'et\'eo (avant fusion) \\
06\_train\_yolo & entra\^inement \\
07\_leak\_check, 08\_novelty\_recall & contr\^ole de fuite train/test et recall selon la nouveaut\'e de la vue \\
09\_fusion\_depth & fusion LiDAR : profondeur, comparaison \`a deux baselines cam\'era, rejet de fausses alarmes \\
10\_aeb\_sweep & balayage des seuils AEB sur un mod\`ele cin\'ematique (sans CARLA) \\
11\_aeb\_closed\_loop & un essai AEB en boucle ferm\'ee dans CARLA \\
12\_aeb\_campaign & campagne : vitesses $\times$ m\'et\'eos $\times$ graines, tableaux agr\'eg\'es \\
13\_agent\_eval & \'evalue la conversion langage naturel $\rightarrow$ JSON (sans CARLA) \\
\bottomrule\end{longtable}

%=====================================================================
\section{Banc de simulation et collecte de donn\'ees}
%=====================================================================
\subsection{Capteurs et param\`etres}
Les param\`etres sont centralis\'es dans \code{src/fusion/config.py} (classe immuable \code{Config}), ce qui permet de
retrouver exactement les conditions d'une collecte.
\begin{center}\begin{tabular}{@{}p{3.6cm}p{11.4cm}@{}}\toprule
Param\`etre & Valeur \\\midrule
Pas de simulation & 0,05~s (20~Hz), mode synchrone \\
Sauvegarde & 1 image sur 4 (5~Hz) \\
Ville de collecte & Town10HD\_Opt \\
\'Episodes & 60, 60 images cibles chacun (12~s de conduite), 70 v\'ehicules de trafic \\
M\'et\'eos & \code{clear} (ClearNoon), \code{rain} (HardRainNoon), \code{fog} (nuages 80, brouillard 60, distance 20~m) \\
Cam\'era RGB & 800$\times$600, FOV 90$^\circ$, position (1,5\,; 0\,; 2,4)~m \\
LiDAR & 64 voies, $1{,}2\cdot10^6$ pts/s, port\'ee 70~m, FOV $+5^\circ/-20^\circ$, \code{dropoff\_general\_rate}=0 \\
\'Etiquetage & distance maximale 60~m ; visibilit\'e minimale 0,3 \\
D\'ecoupage & par \'episode : 70\,\% entra\^inement, 15\,\% validation, 15\,\% test \\
\bottomrule\end{tabular}\end{center}

\begin{choix}{d\'ecoupage par \'episode, pas par image}
Deux images successives d'un m\^eme \'episode sont presque identiques. Un d\'ecoupage al\'eatoire par image mettrait des quasi-doublons
en entra\^inement et en test et gonflerait artificiellement les scores. Chaque \'episode entier est donc affect\'e \`a un seul ensemble,
et les trois m\'et\'eos d'un m\^eme \'episode partagent les m\^emes graines (m\^emes trajectoires) pour comparer proprement les conditions.
Avec 60 \'episodes cela donne 9 \'episodes de test (confirm\'e par les fichiers d'\'evaluation).
\end{choix}

\begin{choix}{LiDAR sans suppression de points}
CARLA supprime 45\,\% des points par d\'efaut (\code{dropoff\_general\_rate}). Pour mesurer la fusion sans ce bruit artificiel non
repr\'esentatif d'un capteur r\'eel particulier, le taux est mis \`a 0. Cons\'equence : le LiDAR simul\'e est plus favorable que du r\'eel,
et il n'est pas d\'egrad\'e par la pluie ni le brouillard (voir limites).
\end{choix}

\subsection{\'Etiquetage automatique}
\code{rig.label\_objects} produit pour chaque v\'ehicule dans le champ :
\begin{itemize}
\item la bo\^ite 2D serr\'ee, obtenue par la \textbf{segmentation d'instance} de CARLA (pas par projection de bo\^ites 3D, qui sur-estime les v\'ehicules partiellement cach\'es) ;
\item la \textbf{visibilit\'e} (pixels visibles / surface de l'enveloppe projet\'ee) ; en dessous de 0,3 le v\'ehicule est marqu\'e \og ignor\'e \fg{} ;
\item les profondeurs de r\'ef\'erence : face avant la plus proche (\code{depth\_face}) et centre ;
\item le nombre de points LiDAR tombant sur l'instance ;
\item le nombre de v\'ehicules dans le champ mais cach\'es (compt\'es, non \'etiquet\'es).
\end{itemize}
\begin{choix}{zone \og ignor\'ee \fg{} plut\^ot que faux n\'egatif}
Un v\'ehicule presque totalement masqu\'e n'est ni un succ\`es ni un \'echec du d\'etecteur. En entra\^inement, son emplacement est gris\'e
dans l'image (les pixels des v\'ehicules gard\'es sont restaur\'es par-dessus) pour ne pas apprendre \`a l'ignorer \`a tort ; en \'evaluation,
une d\'etection sur un v\'ehicule ignor\'e n'est compt\'ee ni comme r\'eussite ni comme fausse alarme (voir \code{05\_eval\_yolo.py}).
\end{choix}

\subsection{Contr\^oles de qualit\'e du jeu de donn\'ees}
\begin{itemize}
\item \code{03\_audit} : audit g\'en\'eral du jeu de donn\'ees ;
\item \code{07\_leak\_check} : pour chaque image de test, distance \`a l'image d'entra\^inement la plus proche de m\^eme cap ($\pm30^\circ$) ;
\item \code{08\_novelty\_recall} : recall du d\'etecteur selon cette distance (familier $<5$~m, interm\'ediaire, nouveau $\geq20$~m).
\end{itemize}
Ces scripts existent et ont \'et\'e utilis\'es, mais leurs sorties chiffr\'ees ne sont pas rejou\'ees dans ce rapport : elles n'ont pas \'et\'e
conserv\'ees dans \code{reference/}. Je ne cite donc pas de pourcentage de fuite.

%=====================================================================
\section{D\'etecteur YOLO}
%=====================================================================
\subsection{Choix et entra\^inement}
\begin{itemize}
\item Mod\`ele de d\'epart : \code{yolo26n.pt} (le plus petit, adapt\'e \`a un GPU de 6~Go et \`a une inf\'erence en boucle ferm\'ee \`a 20~Hz).
\item Une seule classe \emph{vehicle} : la perception sert \`a la distance d'un obstacle, pas \`a la cat\'egorie.
\item Param\`etres (\code{reference/train\_v4\_n\_args.yaml}) : 40 \'epoques maximum, patience 10, lot de 8, images 800 px, optimiseur automatique,
      mosa\"ique, retournement horizontal, HSV, graine 0, d\'eterministe.
\item L'entra\^inement s'est arr\^et\'e \`a l'\'epoque 38 (patience), meilleur mod\`ele \`a l'\'epoque 28, soit environ 76~min.
\end{itemize}
Meilleure \'epoque sur la validation : mAP50 = 0,948 ; mAP50--95 = 0,771 ; pr\'ecision 0,908 ; recall 0,900.
L'empreinte SHA-256 du fichier de poids est conserv\'ee dans \code{reference/best\_pt.sha256} pour retrouver exactement le mod\`ele utilis\'e.

\subsection{R\'esultats sur le jeu de test (recall par distance et m\'et\'eo)}
Recall = part des v\'ehicules \'etiquet\'es (non ignor\'es) retrouv\'es par une d\'etection, par tranche de distance. Les chiffres ci-dessous
sont recalcul\'es depuis les CSV de \code{reference/}.
\begin{center}\small\begin{tabular}{@{}llccc@{}}\toprule
Mod\`ele, ville & M\'et\'eo & 0--20 m & 20--40 m & 40--60 m \\\midrule
YOLO26n pr\'e-entra\^in\'e (COCO)$^{*}$, Town10 & clair & 99,6\,\% & 91,9\,\% & 69,7\,\% \\
 & pluie & 98,5\,\% & 78,5\,\% & 44,0\,\% \\
 & brouillard & 96,4\,\% & 48,2\,\% & 3,4\,\% \\\midrule
Fine-tun\'e, Town10 (4\,451 objets, 9 \'ep.) & clair & 100\,\% & 98,8\,\% & 89,3\,\% \\
 & pluie & 100\,\% & 99,0\,\% & 85,5\,\% \\
 & brouillard & 100\,\% & 99,0\,\% & 89,7\,\% \\\midrule
Pr\'e-entra\^in\'e (COCO), Town03 (3\,393 objets, 10 \'ep.) & clair & 98,2\,\% & 87,3\,\% & 76,4\,\% \\
 & pluie & 97,2\,\% & 74,3\,\% & 60,3\,\% \\
 & brouillard & 97,5\,\% & 62,0\,\% & 6,9\,\% \\\midrule
Fine-tun\'e, Town03 & clair & 100\,\% & 96,8\,\% & 86,7\,\% \\
 & pluie & 100\,\% & 95,9\,\% & 82,0\,\% \\
 & brouillard & 100\,\% & 95,5\,\% & 89,7\,\% \\
\bottomrule\end{tabular}\end{center}
{\footnotesize $^{*}$ d'apr\`es le nom du fichier (\code{eval\_yolo26n\_test.csv}) et l'usage document\'e de \code{05\_eval\_yolo.py --coco} ; classes COCO voiture, bus, camion.}

\begin{honnete}
Le fine-tuning r\'esout surtout le brouillard \`a moyenne et longue distance (de 3,4\,\% \`a 89,7\,\% \`a 40--60~m). Mais le d\'etecteur est
entra\^in\'e sur \emph{CARLA} Town10 : ces scores mesurent un d\'etecteur de rendu synth\'etique sur des mod\`eles de v\'ehicules
identiques \`a l'entra\^inement, pas une performance sur des images r\'eelles. Town03 est une seconde ville, mais l'absence de fuite
d'entra\^inement sur ses donn\'ees reste \`a confirmer d'apr\`es l'organisation des dossiers (\code{data\textbackslash v4\_town03}) avant toute revendication.
\end{honnete}

%=====================================================================
\section{Fusion LiDAR--cam\'era : distance et rejet de fausses alarmes}
%=====================================================================
\subsection{M\'ethode (fusion tardive)}
La fusion est \textbf{tardive} (au niveau des objets d\'etect\'es) : le d\'etecteur 2D produit des bo\^ites, puis le LiDAR fournit la distance.
\begin{enumerate}
\item Les points LiDAR sont transform\'es dans le rep\`ere cam\'era, puis projet\'es dans l'image (\code{project\_lidar}).
\item Les points du \textbf{sol} sont retir\'es : seuls ceux situ\'es \`a plus de 0,30~m au-dessus du sol (\code{GROUND\_MARGIN}) sont gard\'es.
\item On ne garde que les points dans la \textbf{zone centrale} de la bo\^ite (60\,\% de sa largeur et de sa hauteur, \code{INNER}).
\item La distance est celle du \textbf{premier groupe} de profondeur (le plus proche) : c'est la face avant de l'obstacle, ce qui compte pour freiner.
\item Sans point LiDAR : pas de mesure (utile pour le rejet, et pour l'AEB : \emph{pas de profondeur = pas de piste}).
\end{enumerate}
\begin{choix}{fusion tardive}
Elle exploite un d\'etecteur 2D standard, garde chaque capteur diagnostiquable s\'epar\'ement (on peut mesurer ce que chacun apporte) et co\^ute
peu de calcul. Une fusion pr\'ecoce (nuage de points color\'e) ou au niveau des caract\'eristiques demanderait un r\'eseau d\'edi\'e : hors de port\'ee du GPU
et du temps disponibles, et plus difficile \`a valider en boucle ferm\'ee.
\end{choix}
\begin{choix}{premier groupe plut\^ot que m\'ediane}
La m\'ediane des points de la bo\^ite m\'elange l'obstacle et l'arri\`ere-plan visible entre les roues ou au-dessus du toit. Le premier groupe
donne la face avant, qui est la distance de freinage pertinente.
\end{choix}

\subsection{Comparaison \`a deux estimations cam\'era seule}
R\'ef\'erence : face avant du v\'ehicule (\'etiquette CARLA). Erreur = estim\'e $-$ r\'ef\'erence. Deux baselines cam\'era : (i) \textbf{plan du sol}
(la distance se d\'eduit du bas de la bo\^ite, \code{ground\_depth} $=f_y\,h_{cam}/(y_{bas}-v_0)$) ; (ii) \textbf{hauteur fixe} 1,5~m.
Town10HD\_Opt : 4\,256 d\'etections justes ; Town03 : 3\,172.
\begin{center}\small\begin{tabular}{@{}lccc@{}}\toprule
Ville, distance & LiDAR : MAE (part $<1$~m) & Plan du sol : MAE (part $<1$~m) & Hauteur fixe 1,5~m \\\midrule
Town10, 0--20 m & 0,83 (74\,\%) & 1,01 (73\,\%) & 0,90 (69\,\%) \\
Town10, 20--40 m & 1,20 (75\,\%) & 0,99 (57\,\%) & 3,36 (31\,\%) \\
Town10, 40--60 m & 1,67 (78\,\%) & 1,29 (50\,\%) & 6,27 (13\,\%) \\
Town03, 0--20 m & 0,71 (77\,\%) & 0,98 (73\,\%) & 0,77 (73\,\%) \\
Town03, 20--40 m & 1,10 (81\,\%) & 2,32 (59\,\%) & 3,46 (21\,\%) \\
Town03, 40--60 m & 0,68 (93\,\%) & 7,10 (33\,\%) & 5,43 (12\,\%) \\
\bottomrule\end{tabular}\end{center}

\begin{honnete}
Le LiDAR n'est \emph{pas} plus pr\'ecis en moyenne partout : \`a Town10, le plan du sol a une MAE plus faible \`a 20--60~m (1,0 et 1,3~m contre 1,2 et 1,7~m),
mais seulement 50--57\,\% des d\'etections \`a moins de 1~m contre 75--78\,\%. \`A Town03, le plan du sol se d\'egrade nettement (MAE de 7,1~m \`a 40--60~m ;
cause non v\'erifi\'ee, probablement la pente de la route qui viole l'hypoth\`ese de route plane). Le LiDAR reste stable d'une ville \`a l'autre
(MAE de 0,7 \`a 1,7~m, couverture de 99 \`a 100\,\%, biais m\'edian de +0,2 \`a +0,6~m, corrig\'e du pare-chocs dans l'AEB). La hauteur fixe est mauvaise au-del\`a de 20~m.
\textbf{Il ne faut donc pas affirmer que le LiDAR bat la baseline plan-du-sol en pr\'ecision \`a Town10} ; l'avantage est la robustesse (d'une ville \`a l'autre,
et part de d\'etections $<1$~m).
\end{honnete}

\subsection{Rejet des d\'etections sans preuve LiDAR}
Une d\'etection n'est conserv\'ee que si au moins $N$ points LiDAR tombent dans sa bo\^ite.
\begin{center}\small\begin{tabular}{@{}lcccc@{}}\toprule
 & \multicolumn{2}{c}{Town10} & \multicolumn{2}{c}{Town03} \\
 & pr\'ecision & recall & pr\'ecision & recall \\\midrule
Sans rejet & 92,1\,\% & 95,6\,\% & 86,9\,\% & 93,5\,\% \\
$N=3$ & 92,6\,\% & 95,5\,\% & 90,2\,\% & 92,9\,\% \\
$N=10$ & 94,6\,\% & 86,1\,\% & 94,0\,\% & 79,1\,\% \\
\bottomrule\end{tabular}\end{center}
Le rejet \`a $N=3$ retire 7\,\% des fausses alarmes \`a Town10 et 28\,\% \`a Town03. L'effet est modeste ; $N=10$ co\^ute trop de recall.
Ces fausses alarmes sont celles du d\'etecteur sur des images CARLA, pas une mesure sur du trafic r\'eel. Le seuil retenu pour l'AEB exige simplement une
mesure LiDAR (le tracker n'accepte pas de d\'etection sans profondeur).

\subsection{Mesure lat\'erale}
Pour d\'ecider si une cible est dans le couloir de trajectoire, on a besoin de sa position lat\'erale. Elle est calcul\'ee \`a partir des points LiDAR
(option \code{--lateral-source lidar}, d\'efaut) plut\^ot qu'\`a partir du centre de la bo\^ite cam\'era. \teste, mais cette mesure \emph{n'a pas r\'esolu seule}
le freinage parasite \`a 5~Hz (voir section~\ref{sec:bug5hz}).

%=====================================================================
\section{Suivi de cibles (filtre de Kalman)}
%=====================================================================
Le tracker travaille dans le rep\`ere du v\'ehicule ($x$ devant, $y$ lat\'eral, droite positive). Les vitesses sont \textbf{relatives} \`a l'ego :
un objet immobile devant un ego \`a 14~m/s a $v_x=-14$~m/s.
\begin{itemize}
\item \textbf{Mod\`ele \code{cv}} : \'etat $(x,y,v_x,v_y)$, vitesse constante (r\'ef\'erence historique).
\item \textbf{Mod\`ele \code{ca}} (d\'efaut) : \'etat $(x,y,v_x,v_y,a_x,a_y)$, acc\'el\'eration constante avec bruit de saccade (\emph{jerk}, 2~m/s$^3$ dans le script 11).
\item Mesure : distance de la face avant par LiDAR, \'ecart-type de position 0,3~m ; association par porte ($2+25\,\Delta t$ m) ; piste confirm\'ee apr\`es plusieurs mesures ; supprim\'ee apr\`es 3 pertes cons\'ecutives.
\item Initialisation : \'ecart-type de vitesse longitudinale 8~m/s (une piste neuve ignore le rapprochement), mais \textbf{vitesse lat\'erale 2~m/s} (voir \S\ref{sec:bug5hz}).
\end{itemize}
\begin{choix}{mod\`ele \`a acc\'el\'eration constante}
Avec un mod\`ele \`a vitesse constante, une cible qui freine est vue comme \og encore rapide \fg{} pendant plusieurs images. L'acc\'el\'eration relative estim\'ee permet
d'\emph{anticiper} l'arr\^et de la cible (voir AEB). Prix : une estimation plus bruit\'ee, donc un seuil minimal de d\'ec\'el\'eration pour ne pas r\'eagir au bruit.
\end{choix}
\code{Track.misses} compte les images cons\'ecutives sans mesure ; il sert \`a la garde contre les pistes en roue libre (section suivante).

%=====================================================================
\section{AEB : m\'ethode et choix de conception}
%=====================================================================
\subsection{Machine d'\'etats}
\begin{figure}[H]
\centering
\resizebox{\textwidth}{!}{\begin{tikzpicture}[st/.style={draw=accent,rounded corners,fill=accent!6,minimum width=30mm,minimum height=9mm,font=\footnotesize,align=center},
 arr/.style={-{Stealth[length=2mm]},thick,accent},node distance=8mm]
\node[st] (n) {NORMAL};
\node[st,right=of n] (w) {WARNING\\alerte seule};
\node[st,right=of w] (p) {PARTIAL\_BRAKING\\frein 70\,\%};
\node[st,right=of p] (e) {EMERGENCY\_BRAKING\\frein 100\,\%};
\draw[arr] (n)--(w); \draw[arr] (w)--(p); \draw[arr] (p)--(e);
\draw[arr,dashed,grayc] (e.south) -- ++(0,-.9) -| node[pos=.25,below,font=\scriptsize,text=grayc]{descente : hyst\'er\'esis $0{,}7$, attente 1~s, aucun obstacle proche qui ne s'\'eloigne} (n.south);
\end{tikzpicture}}
\caption{\'Etats de l'AEB. L'escalade exige la confirmation du niveau sur plusieurs images ; la descente est volontairement lente.}
\end{figure}

\subsection{Crit\`ere principal : la d\'ec\'el\'eration requise}
Pour chaque cible dans le couloir ($|y|<1{,}5$~m, distance $<60$~m, rapprochement $>0{,}5$~m/s), on calcule la d\'ec\'el\'eration n\'ecessaire pour
s'arr\^eter avant l'obstacle, latence comprise :
\[
 a_{req}=\frac{v_c^{2}}{2\,\bigl(x-v_c\,t_{lat}-m\bigr)},\qquad t_{lat}=0{,}4~\text{s},\; m=1~\text{m},
\]
infinie si le d\'enominateur est n\'egatif ou nul. $v_c$ est la vitesse de rapprochement, $x$ la distance. Niveaux : alerte si $a_{req}\ge 2{,}0$~m/s$^2$ ou
TTC $<2{,}6$~s ; freinage partiel (frein 70\,\%) si $a_{req}\ge 3{,}0$ ; urgence (frein 100\,\%) si $a_{req}\ge 4{,}0$.

\begin{choix}{$a_{req}$ plut\^ot qu'un seuil de TTC fixe}
Un TTC fixe donne la m\^eme alerte \`a 30 et \`a 70~km/h alors que la distance d'arr\^et varie comme $v^2$ : un essai a montr\'e qu'un seuil de TTC freine trop tard
\`a 50~km/h avec un d\'elai d'actionnement de 0,2~s. La d\'ec\'el\'eration requise tient compte de la vitesse, de la distance et de la latence. Le TTC ne sert qu'\`a l'alerte.
\end{choix}
\begin{choix}{seuils calibr\'es sur la voiture CARLA, pas sur une norme}
La premi\`ere version supposait 8~m/s$^2$ de d\'ec\'el\'eration maximale. La mesure donne environ 4,1~m/s$^2$ \`a 70~km/h (frein 100\,\%) pour le v\'ehicule de test.
Les seuils (2,0 / 3,0 / 4,0) sont des valeurs de d\'epart choisies apr\`es balayage en boucle ferm\'ee (jeu \og B \fg) ; ce ne sont pas des valeurs homologu\'ees.
Le jeu pr\'ec\'edent est gard\'e (\code{AEBConfig.conservative()} : 1,5 / 2,5 / 3,5, frein partiel 60\,\%) ; il freine environ 0,7~s plus t\^ot \`a 50~km/h mais
d\'eclenche \`a TTC $\approx3$~s, plus difficile \`a juger comme n\'ecessaire.
\end{choix}

\subsection{Mesures de robustesse}
\begin{enumerate}[label=\textbf{\arabic*.}]
\item \textbf{Confirmation} : le niveau doit \^etre demand\'e sur 2 images cons\'ecutives de l'AEB, et la piste doit persister 0,15~s
      (4 mesures \`a 20~Hz, 3 \`a 10~Hz, 2 \`a 5~Hz) : filtre les fausses d\'etections courtes.
\item \textbf{Hyst\'er\'esis} : un niveau n'est quitt\'e que si $a_{req}$ tombe sous $0{,}7\times$ son seuil. Sans cela, un freinage qui r\'eussit fait osciller $a_{req}$ autour du seuil,
      le frein est rel\^ach\'e \`a 16~m d'un obstacle fixe, et il faut ensuite freiner plus fort (observ\'e en CARLA \`a 50~km/h).
\item \textbf{Maintien} : on ne rel\^ache pas tant qu'un obstacle du couloir est \`a moins de 15~m sans s'\'eloigner ; un essai a montr\'e qu'un rel\^achement \`a 1~m/s avec un conducteur
      qui r\'e-acc\'el\`ere ram\`ene l'ego \`a 0,5~m de l'obstacle.
\item \textbf{Anticipation de la d\'ec\'el\'eration de la cible} (\code{--lead-decel}, activ\'ee dans le script 11 avec le tracker \code{ca}) : si l'acc\'el\'eration absolue estim\'ee de la cible est
      $<-3$~m/s$^2$ sur 3 images, on suppose qu'elle s'arr\^ete et on calcule la d\'ec\'el\'eration pour s'arr\^eter avant son point d'arr\^et (plafond 8~m/s$^2$).
\item \textbf{Latence tenant compte de la fr\'equence de perception} : quand la perception tourne \`a moins de 20~Hz, l'AEB ajoute environ 3 p\'eriodes de calcul \`a la latence
      (maintien du frein + confirmation) ; c'est une hypoth\`ese de mod\`ele.
\item \textbf{Garde contre les pistes en roue libre} (\code{escalate\_max\_misses}=2) : une piste sans mesure depuis plus de 2 images est extrapol\'ee, non observ\'ee ;
      elle ne peut plus faire \emph{monter} le niveau, mais retient un frein d\'ej\`a engag\'e. Une tol\'erance de 0 a \'et\'e \'ecart\'ee : jusqu'\`a +9 collisions sur 40 \`a 5~Hz avec 30\,\% de pertes (banc cin\'ematique).
\end{enumerate}

%=====================================================================
\section{Validation en boucle ferm\'ee dans CARLA}
%=====================================================================
\subsection{Banc de test}
Le script \code{11\_aeb\_closed\_loop} lance un sc\'enario avec seulement l'ego et un v\'ehicule cible (\textbf{sans trafic}). \`A chaque tick de 20~Hz :
capteurs $\rightarrow$ YOLO $\rightarrow$ LiDAR $\rightarrow$ tracker $\rightarrow$ AEB $\rightarrow$ commande de l'ego. Le conducteur et le r\'egulateur de l'ego sont \'ecrits pour l'essai.
Sorties par run : \code{log.csv} (une ligne par tick), \code{summary.json}, images annot\'ees, \code{status.json}.

\begin{center}\small\begin{tabular}{@{}lp{11cm}@{}}\toprule
Sc\'enario & Description et attente \\\midrule
\code{stopped} & cible \`a l'arr\^et dans la voie ; l'AEB doit \'eviter la collision \\
\code{braking} & cible qui roule puis freine \`a $t=$ \code{--lead-brake-t} (frein 0,75 par d\'efaut) \\
\code{adjacent} & v\'ehicule \`a l'arr\^et en voie voisine (d\'ecalage 3,5~m) : l'ego doit passer \textbf{sans freiner} \\
\code{follow} & cible \`a vitesse constante, sans freiner : l'ego ne doit \textbf{pas} freiner \\
\code{approach} & approche avec un r\'egulateur de distance d'un v\'ehicule plus lent : conduite normale, pas de freinage \\
\bottomrule\end{tabular}\end{center}
Les sc\'enarios \emph{adjacent}, \emph{follow} et \emph{approach} mesurent les \textbf{freinages parasites} ; les deux premiers mesurent l'\textbf{efficacit\'e}.
Les campagnes (\code{12\_aeb\_campaign}) r\'ep\`etent un cas sur des vitesses, des m\'et\'eos et des graines ; \`a chaque essai un processus CARLA distinct est lanc\'e (\,$\approx15$~s).

\subsection{R\'esultats, Town10HD\_Opt, 20~Hz, m\'et\'eo claire}
\begin{center}\small\begin{tabular}{@{}p{7.3cm}p{8cm}@{}}\toprule
Sc\'enario & R\'esultat \\\midrule
Cible arr\^et\'ee, 30 / 50 / 70~km/h (3 graines, \'ecart de d\'epart 50~m) & 0 collision ; \'ecart minimal m\'edian 9,9 / 7,6 / 3,3~m ; TTC au premier freinage 1,8 / 2,5 / 2,4~s \\
Idem avec le jeu conservateur & 0 collision ; 11,9 / 9,2 / 4,3~m ; TTC 2,1 / 3,2 / 2,4~s \\
Cible qui freine (1~s apr\`es le d\'epart), 50~km/h, \'ecart 20~m (2 graines) & 0 collision, \'ecart minimal 8,6~m (\emph{sans} anticipation : 2/2 collisions, impact 23~km/h) \\
70~km/h avec plus de place (d\'epart 60~m, 3 graines) & 0 collision, \'ecart minimal m\'edian 13,2~m \\
Cible hors couloir, suivi \`a m\^eme vitesse, approche normale & 0 freinage parasite observ\'e (alerte seule en approche \`a 20 et 18~m) \\
\bottomrule\end{tabular}\end{center}
\begin{honnete}
\`A 70~km/h avec 50~m de d\'epart, la distance est d\'ej\`a inf\'erieure \`a la distance d'arr\^et : le freinage part d\`es la premi\`ere d\'etection quel que soit le seuil.
Cette ligne ne d\'epartage donc pas les jeux de seuils. Les essais ont 2 \`a 4 graines : ils montrent un comportement, pas une distribution statistique.
\end{honnete}

\subsection{Effet de la fr\'equence de perception}
Cible arr\^et\'ee, d\'epart 60~m, 2 graines, \'ecart minimal m\'edian :
\begin{center}\small\begin{tabular}{@{}lccc@{}}\toprule
Fr\'equence & 20 Hz & 10 Hz & 5 Hz \\\midrule
50 km/h & 7,6 m (d\'epart 50 m) & 8,0 m & 5,3 m \\
70 km/h & 13,2 m & 11,5 m & 6,5 m \\
\bottomrule\end{tabular}\end{center}
Sans correction de la latence suppos\'ee, \`a 5~Hz le m\^eme jeu provoquait une collision (13~km/h) sur 1 essai sur 2 \`a 50~km/h. Les effets de cette correction sont valid\'es sur la cible arr\^et\'ee
seulement, pas sur les sc\'enarios non dangereux \`a 5 et 10~Hz.

\subsection{Seconde ville : Town03}
Mêmes scripts et r\'eglages, 2 graines (42 et 43), CARLA lanc\'e \textbf{sans} l'option \code{-quality-level=Low}.
\begin{limite}[Cause non \'elucid\'ee]
Avec \code{-quality-level=Low}, le chargement de Town03 plante (\og Shader compilation failures are Fatal \fg). Le rendu par d\'efaut d\'eplace les \'ecarts \`a 50~km/h de moins d'un m\`etre :
le t\'emoin Town10 a donc \'et\'e refait dans le m\^eme mode, et on ne m\'elange pas deux qualit\'es dans une comparaison.
\end{limite}
\begin{center}\small\begin{tabular}{@{}p{5.4cm}p{3.2cm}p{3cm}p{3cm}@{}}\toprule
Sc\'enario (\'ecart 50~m sauf mention) & Town03, d\'efaut & Town10, m\^eme rendu & Town10, Low \\\midrule
Cible arr\^et\'ee 30/50/70 : collisions & 0/2 chacun & 0/2 chacun & 0/2 chacun \\
idem : \'ecart minimal m\'edian (m) & 9,9 / 9,6 / 3,9 & 10,2 / 8,4 / 3,3 & 9,9 / 7,7 / 3,3 \\
Cible qui freine (m\^eme vitesse), \'ecart 20 m, 50 km/h & 0/2, 8,5 m & 0/2, 8,6 m & 0/2, 8,6 m \\
Cible qui freine (30 km/h), \'ecart 45 m, 50 km/h & 0/2, 7,5 m (pire 6,7) & 0/2, 8,0 m & non mesur\'e \\
Hors couloir 30/50/70, 20 et 5 Hz & 0 freinage, 0 alerte (3 graines) & non refait & 0 freinage (6 graines, 5 Hz) \\
Suivi m\^eme vitesse (15 m), 30/50 km/h & 0/4 freinages & non refait & 0/4 \\
\bottomrule\end{tabular}\end{center}
Les \'ecarts entre villes sont du m\^eme ordre ($\le1{,}2$~m pour des \'ecarts de 7 \`a 10~m, avec 2 graines) : ce sont des indications, pas une diff\'erence mesur\'ee.

\subsection{Un cas non \'evitable, volontairement exclu des tableaux}
Une cible \`a 30~km/h qui s'arr\^ete en 0,5~s avec 20~m d'\'ecart initial (ego \`a 50~km/h) donne 2 collisions sur 2 (impact $\approx29$~km/h) \`a Town03 comme \`a Town10. Le journal montre que l'alerte
arrive \`a 0,2~s, mais que l'arr\^et de la cible survient \`a 12--13~m ; la d\'ec\'el\'eration n\'ecessaire ($\approx4{,}8$~m/s$^2$, calcul approch\'e) d\'epasse celle de la voiture ($\approx4{,}1$ mesur\'ee),
m\^eme sans latence. Un AEB ne corrige pas une situation physiquement inatteignable. Une erreur de ma part, rep\'er\'ee ici : une variante (\code{--gap 20} avec une cible \`a 30~km/h)
avait \'et\'e propos\'ee sans validation ; elle est remplac\'ee par des variantes valid\'ees (m\^eme vitesse, \'ecart 20~m\,; cible \`a 30~km/h, \'ecart 45~m).

%=====================================================================
\section{Robustesse aux fausses d\'etections}
%=====================================================================
\subsection{Injection de d\'efauts}
Le module \code{FaultInjector} permet de simuler : des fausses d\'etections dans le couloir (probabilit\'e par image, dur\'ee r\'eglable), des pertes de mesure, un bruit gaussien sur la distance.
C'est un \textbf{test de contrainte}, pas un taux r\'ealiste : les statistiques de persistance des fausses d\'etections du d\'etecteur r\'eel ne sont pas mesur\'ees.

\subsection{M\'ethode : pr\'edire puis mesurer}
Les fr\'equences d'essais avec au moins un freinage parasite ont d'abord \'et\'e \emph{pr\'edites} par le mod\`ele hors CARLA (\code{tests/test\_faults.py}, 200 simulations), puis \emph{mesur\'ees} en boucle ferm\'ee
(50~km/h, route d\'egag\'ee, probabilit\'e 0,10 par image, $\approx8$ \'ev\'enements par essai de 4~s, 6 graines).
\begin{center}\small\begin{tabular}{@{}lccc@{}}\toprule
Dur\'ee & Mod\`ele hors CARLA & CARLA, confirm. 2 mesures & CARLA, confirm. 0,15~s \\\midrule
1 image & 14\,\% & 1/6 & 0/6 \\
2 images & 49\,\% & 3/6 & 0/6 \\
3 images & 98,5\,\% & 6/6 & 0/6 \\
\bottomrule\end{tabular}\end{center}
Avec la confirmation de 0,15~s, le mod\`ele pr\'evoit 0\,\%, 17\,\% et 20\,\% pour 1, 2 et 3 images : le r\'esultat 0/6 est \emph{compatible} avec lui, il ne le prouve pas. Une fausse d\'etection de 4 images ou plus
freine toujours (100\,\% pr\'edit, non mesur\'e en CARLA). Co\^ut de la confirmation : environ 0,7~m de retard \`a 50~km/h si l'obstacle appara\^it \`a 20--25~m (mod\`ele hors CARLA) ; pas de r\'egression sur la cible
arr\^et\'ee \`a 20~Hz (m\^emes \'ecarts minimaux et distances de premier freinage).

\subsection{Un bug trouv\'e puis corrig\'e : freinage parasite \`a 5~Hz}\label{sec:bug5hz}
\textbf{Sympt\^ome} : \`a 5~Hz, en passant \`a c\^ot\'e d'une cible arr\^et\'ee dans la voie voisine (50~km/h, graine 43), un freinage d'urgence parasite apparaissait dans 4 lancements sur 6.\\
\textbf{Hypoth\`ese \'ecart\'ee} : l'erreur de mesure lat\'erale. Passer \`a une mesure lat\'erale par points LiDAR n'a pas r\'esolu le probl\`eme \`a lui seul.\\
\textbf{Cause trouv\'ee} en rejouant le journal : une fausse piste n\'ee de deux d\'etections proches recevait une vitesse lat\'erale initiale d'environ $-5$~m/s (a priori trop large de 8~m/s) et d\'erivait dans le couloir pendant 3 images.\\
\textbf{Correction} : \`a l'initialisation, \'ecart-type de la vitesse lat\'erale ramen\'e \`a 2~m/s (\code{vy\_init\_std}). Test de non-r\'egression \code{test\_depassement\_a\_5hz\_ne\_declenche\_pas\_le\_freinage}.\\
\textbf{Validation} : 0 freinage sur 8 lancements de la graine 43, et 0/18 sur les graines 42 \`a 47 \`a 30 / 50 / 70~km/h.\\
\textbf{Garde ajout\'ee ensuite} (\code{escalate\_max\_misses}=2) : neutralise la s\'equence rejou\'ee m\^eme avec l'ancien a priori. Revalid\'ee en CARLA : cible arr\^et\'ee et v\'ehicule qui freine inchang\'es, 0/9 freinage hors couloir \`a 5~Hz.

\begin{choix}{rejouer le journal avant de modifier}
La cause n'\'etait pas la premi\`ere hypoth\`ese. Rejouer les mesures enregistr\'ees dans le tracker a permis de voir la vitesse initiale aberrante, au lieu d'ajuster des seuils \`a l'aveugle
(ce qui aurait masqu\'e le d\'efaut en d\'egradant l'efficacit\'e).
\end{choix}

%=====================================================================
\section{ADAS Validation Studio (application web)}
%=====================================================================
\subsection{Architecture}
\begin{figure}[H]
\centering
\resizebox{\textwidth}{!}{\begin{tikzpicture}[node distance=9mm and 12mm,
  box/.style={draw=accent,rounded corners,fill=accent!6,minimum height=9mm,align=center,font=\footnotesize},
  arr/.style={-{Stealth[length=2mm]},thick,accent}]
\node[box] (ui) {Page web\\\code{index.html}};
\node[box,right=of ui] (api) {API FastAPI\\\code{backend/app/main.py}};
\node[box,right=of api] (scr) {Sous-processus\\scripts 11 / 12};
\node[box,right=of scr] (carla) {Serveur CARLA\\\code{localhost:2000}};
\node[box,below=of api] (llm) {Agent\\(LLM ou r\`egles)};
\node[box,below=of scr] (fs) {\code{runs/}\\log, status, summary};
\draw[arr] (ui)--(api); \draw[arr] (api)--(scr); \draw[arr] (scr)--(carla);
\draw[arr] (api)--(llm); \draw[arr] (scr)--(fs); \draw[arr] (api)--(fs);
\end{tikzpicture}}
\caption{Le serveur ne contient pas de logique de simulation : il lance les scripts en sous-processus et lit les fichiers qu'ils \'ecrivent.}
\end{figure}

\begin{choix}{sous-processus plut\^ot que simulation dans le serveur}
CARLA, YOLO et le GPU sont des ressources lourdes et fragiles ; un plantage de CARLA ne doit pas tuer l'API. Un processus par essai donne une isolation naturelle, un arr\^et propre (fichier \code{STOP}),
et r\'eutilise exactement les scripts valid\'es en ligne de commande. Contrepartie : environ 15~s d'initialisation par essai, et un seul run ou une seule campagne \`a la fois (une instance CARLA, un GPU).
\end{choix}

\subsection{Fonctions}
\begin{itemize}
\item \textbf{Run d'observation} : trafic, YOLO + distances LiDAR en direct.
\item \textbf{Run AEB} en boucle ferm\'ee (script 11) : vue vid\'eo en direct (images \`a 10~Hz avec petite m\'emoire tampon), \'etat de l'AEB, d\'ec\'el\'eration requise, \'ecart ; bilan et courbes \`a la fin.
\item Options : AEB d\'esactiv\'e (t\'emoin), fr\'equence de perception 20/10/5~Hz, injection de d\'efauts.
\item \textbf{Campagnes} (\code{POST /api/campaigns}) : vitesses $\times$ m\'et\'eos $\times$ graines, maximum 60 essais, progression, tableau agr\'eg\'e (collisions, impact moyen, \'ecart m\'edian et pire, \'ecart et TTC au premier freinage),
      reprise sur le dossier de campagne, arr\^et.
\item \textbf{Valeurs par d\'efaut par sc\'enario} et indication de la \textbf{route droite n\'ecessaire}, car le tron\c con droit disponible est limit\'e (95~m \`a Town10HD\_Opt, 80~m \`a Town03).
\end{itemize}

\subsection{Choix et incidents notables}
\begin{itemize}
\item \textbf{Longueur de route} : une erreur \og plus long tron\c con droit : 95~m, il en faut 110 \fg{} a conduit \`a calculer la route requise par sc\'enario
      (freinage : $\text{\'ecart}+5+v_l\,t_b+v_l^2/8+10$ ; suivi/approche : $\text{\'ecart}+5+4v_l+15$ ; adjacent et arr\^et : $\text{\'ecart}+30$), \`a l'afficher dans le formulaire et \`a refuser les cas infaisables avant de lancer CARLA.
\item \textbf{Verrous de fichiers Windows} : le navigateur lisant \code{status.json} provoquait un \code{PermissionError} lors du renommage atomique. Solution : \code{replace\_retry} (r\'eessais de quelques ms) ;
      les \'ecritures de fichiers d'affichage sont non critiques (on saute une mise \`a jour), celles de l'\'etat final sont critiques.
\item \textbf{Chargement de Town03} : d\'elai de \code{load\_world} port\'e \`a 180~s ; la ville pass\'ee en argument du serveur \'etait ignor\'ee, on la charge par l'API.
\item \textbf{Test qui lan\c cait un vrai processus} : un test de refus de campagne lan\c cait une vraie campagne quand CARLA \'etait allum\'e. Il remplace d\'esormais l'\'etat de CARLA et interdit \code{Popen}.
\end{itemize}

\begin{limite}[Statut de l'interface]
L'API des campagnes et la page sont \impl{} et \teste{} sans CARLA. Une campagne \emph{lanc\'ee depuis la page} jusqu'\`a un tableau agr\'eg\'e coh\'erent avec la ligne de commande
(arr\^et\'ee 30/50/70, graines 42 et 43, attendu $\approx$ 10 / 7,6 / 3,3~m, 0 collision) \textbf{n'est pas encore confirm\'ee par l'auteur} : elle n'est donc pas \valide.
\end{limite}

%=====================================================================
\section{Agent : langage naturel $\rightarrow$ sc\'enario d'essai}
%=====================================================================
\subsection{But et principe}
Convertir une phrase comme \og Une voiture \`a l'arr\^et \`a 60~m, j'arrive \`a 50~km/h sous la pluie, essaie les graines 42 et 43 \fg{} en un sc\'enario JSON que le Studio sait lancer, ou la \textbf{refuser avec une raison}.
Le LLM ne lance rien : il \'ecrit un JSON, qui est ensuite \textbf{valid\'e par du code d\'eterministe}.

\subsection{Sch\'ema du sc\'enario}
\begin{center}\small\begin{tabular}{@{}ll@{}}\toprule
Champ & Domaine \\\midrule
\code{scenario} & stopped, braking, adjacent, follow, approach \\
\code{ego\_speed\_kmh} & [10, 90] \\
\code{gap\_m} & [10, 80] \\
\code{lead\_speed\_kmh} & [0, 90] ($\ge5$ pour braking, follow, approach) \\
\code{lead\_brake\_t\_s} & [0, 20] \\
\code{weather}, \code{town}, \code{seeds} & clear/rain/fog ; villes du Studio ; liste de graines \\
\code{perception\_hz} & 20, 10 ou 5 \\
\code{aeb} & bool\'een \\
\code{faults} & ghost\_prob [0\,;0,2], ghost\_len [1\,;30], drop\_prob [0\,;0,9], noise\_sigma [0\,;2] \\
\bottomrule\end{tabular}\end{center}
Chaque sc\'enario a des valeurs par d\'efaut (par ex. \code{stopped} : \'ecart 50~m ; \code{braking} : \'ecart 45~m, cible 30~km/h, freinage \`a 1~s ; \code{follow} : \'ecart 15~m, cible \`a la vitesse de l'ego).

\begin{choix}{trois niveaux de validation apr\`es le LLM}
(1) \textbf{format} : champs connus, types, unit\'es, plages ; (2) \textbf{faisabilit\'e} : la route droite requise est compar\'ee \`a la route disponible dans la ville choisie ;
(3) \textbf{reconstruction} d'un corps de requ\^ete pour l'API (\code{to\_studio\_payload}). Un mod\`ele peut se tromper de fa\c con plausible ; la v\'erification doit donc \^etre ind\'ependante du mod\`ele.
Un refus est un r\'esultat valide : \emph{il vaut mieux refuser qu'ex\'ecuter un essai physiquement impossible}.
\end{choix}

\subsection{Convertisseurs compar\'es}
\begin{itemize}
\item \textbf{Baseline \`a r\`egles} (expressions r\'eguli\`eres, sans LLM) : reproductible, sert de plancher.
\item \textbf{LLM} via un client pour API \og compatibles OpenAI \fg{} (biblioth\`eque standard uniquement ; Groq, OpenRouter, Ollama, Gemini ; r\'eessais sur 429/5xx ; \code{User-Agent} explicite car un pare-feu renvoyait une erreur 403 \og 1010 \fg{} sans en-t\^ete).
      Temp\'erature 0, suppression des blocs \code{<think>}.
\item Deux prompts : v1 (886 caract\`eres) et v2 (1\,300 caract\`eres, qui pr\'ecise que la vitesse donn\'ee est celle de l'ego).
\end{itemize}
Choix du fournisseur : une API gratuite (d\'ecision d'octobre 2026, apr\`es avoir envisag\'e d'abord le mod\`ele le plus puissant). Les limites gratuites de Groq relev\'ees (mod\`eles \code{openai/gpt-oss-120b}, \code{qwen/qwen3.8-27b} ; 30 requ\^etes/min, 1\,000/jour,
8\,000 tokens/min, 200\,000 tokens/jour) proviennent de pages d'agr\'egateurs dat\'ees de mi-octobre 2026, \emph{pas d'une documentation officielle} : \`a rev\'erifier.

\subsection{Jeu A (35 demandes) : r\'esultats}
Jeu \code{reference/agent\_benchmark.jsonl} : 35 demandes en fran\c cais et en anglais, \textbf{dont 7 \`a refuser} (infaisables ou trop vagues), \'ecrit par moi. Un cas est correct si tous les champs coincident
(tol\'erance de 0,5 pour les valeurs $\ge5$) ou si le refus attendu est produit.
\begin{center}\small\begin{tabular}{@{}lc@{}}\toprule
Convertisseur & Cas corrects sur 35 \\\midrule
R\`egles (mots-cl\'es) & 23 (65,7\,\%), reproductible \\
\code{openai/gpt-oss-120b} (Groq), prompt v1 & 33, puis 32 (deux passages) \\
Idem, prompt v2 & 34, puis 33, 33, 32 (moyenne 32,7) \\
\bottomrule\end{tabular}\end{center}
\begin{honnete}
Sur les sc\'enarios que le mod\`ele accepte, les champs sont exacts (aucune erreur d'unit\'e ni de valeur observ\'ee). Les \'echecs sont des \textbf{refus \`a tort} de demandes valides (4 cas selon le prompt)
et, une fois, une demande vague accept\'ee. \`A temp\'erature 0, deux passages identiques ne donnent pas le m\^eme score : une diff\'erence de 1 \`a 2 cas sur 35 n'est pas une am\'elioration.
Le prompt v2 corrige deux cas que j'avais formul\'es de fa\c con ambigu\"e et en casse deux de suivi : \textbf{aucun gain d\'emontr\'e} sur v1.
Surtout, le prompt v2 a \'et\'e ajust\'e d'apr\`es les erreurs de ce jeu : ce n'est pas une mesure sur donn\'ees ind\'ependantes.
\end{honnete}

\subsection{Jeu B (26 demandes) : mesure \`a prompt fig\'e}
Pour obtenir une mesure plus honn\^ete, un second jeu \code{reference/agent\_benchmark\_b.jsonl} a \'et\'e \'ecrit \emph{apr\`es} le gel du prompt v2 : 26 demandes (cat\'egories : familier, multi-param\`etres, unit\'es, d\'efauts, anglais,
infaisable, ambigu), dont 7 \`a refuser. Les r\'eponses de r\'ef\'erence sont valid\'ees par le code (aucune r\'ef\'erence invalide ou infaisable) ; l'oracle obtient 26/26 et la baseline \`a r\`egles 9/26 (35\,\%).
\textbf{Le score du LLM sur ce jeu n'est pas encore mesur\'e} (commande pr\'epar\'ee : \code{13\_agent\_eval.py --bench reference\textbackslash agent\_benchmark\_b.jsonl --repeat 3}, sans retoucher le prompt).
Limite persistante : ce jeu est encore \'ecrit par moi ; un jeu C avec les formulations de l'auteur du projet est pr\'evu.

\subsection{Int\'egration au Studio}
\impl{} et \teste{} sans CARLA ni vrai LLM (tests avec un faux mod\`ele) :
\begin{itemize}
\item \code{GET /api/agent} : mode (\code{llm} ou \code{rules}), fournisseur, mod\`ele. Le LLM est utilis\'e si la variable d'environnement \code{LLM\_API\_KEY} existe c\^ot\'e \textbf{serveur} (ou si le fournisseur est Ollama, local) ; sinon repli sur les r\`egles, avec mention dans l'interface.
\item \code{POST /api/agent/parse} : renvoie soit le sc\'enario normalis\'e et la route n\'ecessaire, soit un refus de type \code{rejet} (avec la raison donn\'ee par le mod\`ele), \code{invalide}, \code{infaisable} ou \code{json\_invalide}. Une erreur r\'eseau ou de quota renvoie 502. \textbf{Ne lance jamais de processus.}
\item Interface : zone de texte, bouton \og Convertir \fg, r\'esultat lisible, bouton \og Appliquer au formulaire de campagne \fg{} ; l'utilisateur relit puis lance lui-m\^eme.
\item La campagne du Studio accepte d\'esormais la carte (\code{town}), l'instant de freinage de la cible, la dur\'ee des fausses d\'etections, les pertes et le bruit.
\end{itemize}
\begin{choix}{l'agent propose, l'humain d\'ecide}
Un LLM peut se tromper de fa\c con plausible. Le formulaire pr\'e-rempli est \textbf{relu et lanc\'e par l'utilisateur}. La cl\'e d'API n'est jamais lue depuis le navigateur ni \'ecrite dans un fichier du d\'ep\^ot :
elle reste dans la variable d'environnement du terminal. (Deux cl\'es ont \'et\'e coll\'ees par erreur dans la conversation : elles doivent \^etre r\'evoqu\'ees.)
\end{choix}
\begin{limite}[Ce qui reste \`a faire sur l'agent]
Pas de RAG, pas de LoRA, pas de boucle d'am\'elioration continue avec regroupement des erreurs en production ; un seul mod\`ele test\'e ; l'agent connaît un seul interlocuteur (sc\'enarios du Studio) ; le mode \og sans AEB \fg{} d'une campagne
n'est pas g\'er\'e directement (le Studio lance l'AEB, et les t\'emoins si la case est coch\'ee). Ces \'el\'ements sont \plan.
\end{limite}

%=====================================================================
\section{Tests automatiques}
%=====================================================================
64 tests passent hors CARLA sur ma copie de travail ; ils doivent \^etre relanc\'es par l'auteur sur sa machine (\code{python -m pytest -q tests}). R\'epartition par fichier :
\begin{center}\small\begin{tabular}{@{}lp{11.2cm}@{}}\toprule
Fichier & Contenu \\\midrule
\code{test\_aeb.py} & machine d'\'etats, d\'ec\'el\'eration requise, hyst\'er\'esis, anticipation, non-r\'egression du freinage parasite \`a 5~Hz, garde en roue libre \\
\code{test\_faults.py} & mod\`ele hors CARLA des fausses d\'etections (pr\'ediction des fr\'equences de freinage parasite) \\
\code{test\_perception.py} & fusion, \code{ground\_depth}, rejet \\
\code{test\_fusion\_lateral.py} & mesure lat\'erale par points LiDAR \\
\code{test\_api.py} & runs, campagnes (commande construite, refus, suivi), agent (LLM simul\'e, refus, JSON illisible, infaisable, repli r\`egles, aucun \code{Popen}) \\
\code{test\_agent.py} & sch\'ema, validation, faisabilit\'e, m\'etriques, baseline, client HTTP (faux serveur local), jeu B, prompts \\
\bottomrule\end{tabular}\end{center}
Les tests v\'erifient le \emph{code}, pas la physique : la physique se v\'erifie en CARLA (sections pr\'ec\'edentes).

%=====================================================================
\section{Registre des difficult\'es rencontr\'ees}
%=====================================================================
\begin{longtable}{@{}p{4.2cm}p{4.8cm}p{5.8cm}@{}}\toprule
Probl\`eme & Cause & Solution / enseignement \\\midrule\endhead
Freinage parasite \`a 5~Hz (graine 43) & a priori de vitesse lat\'erale initiale trop large (8~m/s) & \'ecart-type 2~m/s ; garde \code{escalate\_max\_misses}=2 ; test de non-r\'egression \\
Garde de tol\'erance 0 & les pistes de 1 image sans mesure sont l\'egitimes \`a 5~Hz avec pertes & tol\'erance 2 ; 0 donnait +9 collisions sur 40 \\
Fausses d\'etections de 1--3 images & confirmation trop courte & confirmation 0,15~s ; mod\`ele hors CARLA pr\'edit puis CARLA mesure \\
Rel\^achement du frein trop t\^ot & $a_{req}$ oscille autour du seuil & hyst\'er\'esis 0,7 ; maintien proche d'un obstacle \\
Collision \`a 5~Hz sans correction & latence r\'eelle sup\'erieure \`a celle suppos\'ee & latence tenant compte de la fr\'equence ; valid\'ee sur cible arr\^et\'ee seulement \\
\code{PermissionError} sur \code{status.json} & lecteur (navigateur) ouvre le fichier pendant le renommage & \code{replace\_retry} ; fichiers non critiques saut\'es \\
Route trop courte (110~m requis, 95~m dispos) & scénario choisi sans calcul pr\'ealable & formules de longueur, indication dans l'UI, refus avant lancement \\
Plantage de CARLA \`a Town03 & \code{-quality-level=Low} : \og shader compilation \fg & CARLA sans option ; cause non \'elucid\'ee ; ne pas m\'elanger les qualit\'es \\
Test qui lan\c cait une vraie campagne & CARLA joignable lors des tests & \'etat CARLA simul\'e, garde sur \code{Popen} \\
Cas \og cible 30~km/h, \'ecart 20~m \fg & physiquement inatteignable (4,8 > 4,1~m/s$^2$) & exclu des tableaux ; variantes valid\'ees \\
Erreur HTTP 403 \og 1010 \fg{} (Groq) & absence de \code{User-Agent} & en-t\^ete explicite ; test ajout\'e \\
Cl\'es d'API coll\'ees dans le chat & erreur d'usage & r\'evocation ; variable d'environnement du terminal uniquement \\
Fichiers non \'ecrits sur le disque via le pont & \'ecriture silencieuse sans effet & v\'erification syst\'ematique des tailles apr\`es \'ecriture \\
Scores LLM non reproductibles & temp\'erature 0 non d\'eterministe c\^ot\'e API & r\'ep\'etitions, moyenne, pas de conclusion sur 1--2 cas \\
\bottomrule\end{longtable}

%=====================================================================
\section{Limites et affirmations \`a ne pas faire}
%=====================================================================
\begin{limite}[Limites g\'en\'erales]
\begin{itemize}
\item \textbf{Simulation uniquement} ; deux villes ; 2 \`a 6 graines. Ces essais montrent un comportement ; ils ne constituent ni une validation statistique ni une homologation.
\item Le LiDAR CARLA n'est pas d\'egrad\'e par la pluie ni le brouillard : l'axe m\'et\'eo ne teste que la cam\'era.
\item Mod\`eles de v\'ehicules identiques entre entra\^inement et test ; baseline cam\'era seule na\"ive (hauteur fixe 1,5~m) ; pas de trafic dense dans les essais AEB.
\item Banc de test maison : conducteur et r\'egulateur de l'ego \'ecrits pour l'essai.
\item L'AEB est valid\'e \`a 20~Hz. \`A 5~Hz et 70~km/h, la marge est mince : \'ecart minimal $\approx2{,}4$~m sur 16 lancements sans collision, mais \textbf{une collision} (impact 16~km/h) observ\'ee une fois sur la graine 42
      dans un autre lot, non reproduite, cause non \'elucid\'ee (variabilit\'e du simulateur ou marge insuffisante).
\item Fausse d\'etection de 4 images ou plus : l'AEB freine toujours (non mesur\'e en CARLA).
\item Les freinages parasites ne sont mesur\'es que sur trois sc\'enarios simples.
\item L'agent : jeu de test \'ecrit par moi, prompt ajust\'e dessus, un seul mod\`ele LLM.
\end{itemize}
\end{limite}
\textbf{Affirmations \`a ne pas faire} (par exemple dans un CV ou un entretien) :
\begin{itemize}
\item \og Le LiDAR est plus pr\'ecis que la baseline cam\'era \fg{} (faux \`a Town10 en MAE moyenne).
\item \og 0 collision \`a 5~Hz et 70~km/h \fg{} (une collision observ\'ee).
\item \og L'AEB est valid\'e / homologu\'e \fg{} (calibr\'e sur une voiture simul\'ee, seuils de d\'epart).
\item \og Le d\'etecteur g\'en\'eralise \fg{} \`a des images r\'eelles ou \`a une ville dont l'absence de fuite d'entra\^inement n'est pas confirm\'ee.
\item \og L'agent LLM atteint X\,\% \fg{} sans pr\'eciser qu'il s'agit d'un jeu \'ecrit par moi, avec un prompt ajust\'e sur lui ; un score sur le jeu B ou un jeu C ind\'ependant n'existe pas encore.
\end{itemize}
Formulations d\'efendables : \emph{\og Pipeline LiDAR--cam\'era avec AEB valid\'e en boucle ferm\'ee dans CARLA sur deux villes, avec analyse d'erreurs et corrections document\'ees \fg} ; \emph{\og Agent de conversion langage naturel $\rightarrow$ sc\'enario JSON
avec validation d\'eterministe, \'evalu\'e sur un jeu de demandes avec baseline \`a r\`egles \fg}.

%=====================================================================
\section{Perspectives}
%=====================================================================
\begin{tabular}{@{}p{5cm}p{10cm}@{}}\toprule
\plan & Contenu \\\midrule
Mesure du jeu B & lancer le LLM sur 26 demandes, prompt fig\'e, 3 passages \\
Jeu C ind\'ependant & phrases \'ecrites par l'auteur, r\'eponses attendues corrig\'ees par lui \\
Comparaison Ollama local & mod\`ele de 7--8 milliards de param\`etres quantifi\'e (faisable sur 6~Go de VRAM) : justifierait ou non un LoRA / un RAG \\
RAG et LoRA & seulement si le jeu ind\'ependant montre une marge que le prompt ne comble pas \\
Am\'elioration continue & regroupement des erreurs de l'agent (cat\'egories d\'ej\`a produites par les m\'etriques) puis correction mesur\'ee \\
Validation de l'interface & campagne r\'eelle lanc\'ee depuis la page, comparaison \`a la ligne de commande \\
Perception & mesure de la persistance r\'eelle des fausses d\'etections du d\'etecteur ; trafic dense ; LiDAR d\'egrad\'e par la m\'et\'eo \\
AEB & d\'efense contre les fausses d\'etections de 4 images ou plus ; marge \`a 5~Hz et 70~km/h \\
\bottomrule\end{tabular}

%=====================================================================
\appendix
\section{Reproduire}
%=====================================================================
\begin{lstlisting}
python -m venv venv ; venv\Scripts\activate
pip install -r requirements.txt       # + wheel CARLA (non installable depuis PyPI)
# serveur CARLA lance (sans -quality-level=Low pour Town03)
python scripts\01_collect.py
python -m pytest -q tests             # sans CARLA
python scripts\12_aeb_campaign.py --scenario stopped --speeds 30 50 70 --weathers clear --seeds 42 43 44 45
python -m uvicorn app.main:app --app-dir backend --port 8000     # Studio sur http://localhost:8000

# Agent (la cle reste dans le terminal, jamais dans un fichier du depot)
$env:LLM_API_KEY = Read-Host "cle"
$env:AGENT_PROVIDER = "groq"
python scripts\13_agent_eval.py --model openai --provider groq --llm openai/gpt-oss-120b `
   --sleep 6 --prompt v2 --bench reference\agent_benchmark_b.jsonl --repeat 3
\end{lstlisting}
Les r\'esultats de r\'ef\'erence sont dans \code{reference/} (CSV de fusion, d'\'evaluation du d\'etecteur et des campagnes, empreinte du mod\`ele, jeux de l'agent). Variables du serveur pour l'agent :
\code{LLM\_API\_KEY}, \code{AGENT\_PROVIDER}, \code{AGENT\_MODEL}, \code{AGENT\_BASE\_URL}, \code{AGENT\_PROMPT}.

\section{Glossaire}
\begin{description}[leftmargin=3cm,labelwidth=2.8cm,style=nextline]
\item[ADAS] Advanced Driver Assistance Systems, aides \`a la conduite.
\item[AEB] Autonomous Emergency Braking, freinage d'urgence automatique.
\item[TTC] Time To Collision, temps avant collision $=x/v_c$.
\item[$a_{req}$] D\'ec\'el\'eration requise pour s'arr\^eter avant l'obstacle, latence comprise.
\item[Boucle ferm\'ee] La commande de l'ego modifie ce que les capteurs per\c coivent ensuite (par opposition \`a un rejeu de donn\'ees).
\item[SIL] Software-In-the-Loop : le logiciel \`a tester tourne dans un simulateur.
\item[Fusion tardive] Les capteurs sont combin\'es apr\`es la d\'etection d'objets, pas au niveau des donn\'ees brutes.
\item[MAE] Erreur absolue moyenne.
\item[Graine] Entier qui fixe les al\'eas (trajectoire, trafic) pour rendre un essai reproductible.
\item[RAG / LoRA] R\'ecup\'eration de documents ajout\'es au prompt / adaptation l\'eg\`ere des poids d'un mod\`ele ; \plan{} dans ce projet.
\end{description}

\end{document}
